\documentclass[11pt]{amsart}
\usepackage{amsmath,amssymb,amsthm,url,booktabs}
\newtheorem{theorem}{Theorem}\newtheorem{lemma}[theorem]{Lemma}\newtheorem{proposition}[theorem]{Proposition}
\newtheorem{corollary}[theorem]{Corollary}
\theoremstyle{remark}\newtheorem{remark}[theorem]{Remark}
\title{Conserved charges in linearly reinforced graph urns}
\author{Yitzchak Shmalo}
\address{Einstein Institute of Mathematics, The Hebrew University of Jerusalem, Jerusalem, Israel}
\subjclass[2020]{60K35, 60G42, 62L20}
\keywords{P\'olya urn, reinforcement, graph-based competition, martingale, stochastic approximation}
\begin{document}
\begin{abstract}
In the graph-based P\'olya urn with linear reinforcement (WARM with $\alpha=1$) in continuous time, every vector $\sigma$
in the kernel of the unsigned vertex--edge incidence matrix yields an exact martingale $M_\sigma=\sum_e\sigma_e H(N_e-1)$,
where $H$ is the harmonic number. On every finite graph its terminal value has the law of $\sum_e\sigma_e\log W_e$ with
independent standard exponential variables $W_e$. On the even cycle $C_L$ the normalised counts converge to
$(1+u,1-u,\dots)$ with $L\,\mathrm{artanh}(u)$ equal to the terminal value of the charge, so that $u=\tanh(M/L)$ where $M$
is a sum of $L/2$ independent standard logistic variables, and $\sqrt{L}\,\mathrm{artanh}(u)$ converges in law to
the centred normal law with variance $\pi^2/6$. This gives the law of the random limit for the continuous-time, one-vertex-at-a-time analogue of
the urn studied by Chen and Lucas and by Lima. On the integer line the expected quadratic variation of the charge per
vertex is $\pi^2/6$.
\end{abstract}
\maketitle

\section{The model and the result}
In the urn with graph-based competition of van der Hofstad, Holmes, Kuznetsov and Ruszel \cite{HHKR} (WARM), each edge
of a graph carries a count, each vertex carries a Poisson clock, and when a vertex rings it reinforces one of its edges
with probability proportional to $N_e^\alpha$. The discrete-time urn, in which every edge adds a ball to one of its two
endpoints at each step, was introduced by Bena\"im, Benjamini, Chen and Lima \cite{BBCL}. At $\alpha=1$ it was studied
by Chen and Lucas \cite{CL}, who proved convergence on regular bipartite graphs to a random point of a segment of
equilibria and left the law of that point open, and by Lima \cite{Lima}, who completed the convergence proof for
balanced bipartite graphs and asked for the distribution of the limit (his Question~1). On a cycle the two models are
related by exchanging vertices and edges: the bins of the discrete-time urn on $C_L$ correspond to the edges of WARM on
$C_L$, and its competing pairs to the vertices of WARM. The discrete-time urn updates every pair at each step, whereas
WARM updates one vertex at a time at Poisson times. The results below concern WARM; no transfer of the limit law to the
synchronous urn is proved. On the infinite line, Kleptsyn conjectures that the normalised counts of the critical process
converge to $1$ \cite[Section~5]{AletOP}.

The random limit is fixed by conserved charges. For every $\sigma$ in the kernel of the unsigned incidence matrix,
$M_\sigma=\sum_e\sigma_eH(N_e-1)$ is an exact martingale whose quadratic variation never exceeds
$\zeta(2)\sum_e\sigma_e^2$ (Theorem~\ref{thm:mart}); up to a vanishing error it equals $\sum_e\sigma_e\log(N_e/t)$, which
is conserved by the mean-field flow. Its terminal value has an explicit law on every finite graph:
$\mathbb{E}e^{i\theta M_{\sigma,\infty}}=\prod_e\Gamma(1+i\theta\sigma_e)$, so that the charges jointly have the law of
$(\sum_e\sigma_e\log W_e)_\sigma$ with independent standard exponential $W_e$ (Theorem~\ref{thm:law}). On the even cycle
$C_L$ the kernel is one-dimensional, the normalised counts converge to $(1+u_\infty,1-u_\infty,\dots)$ with
$L\,\mathrm{artanh}(u_\infty)=M_\infty$ (Theorem~\ref{thm:conv}), and therefore $u_\infty=\tanh(M_\infty/L)$ with
$M_\infty$ a sum of $L/2$ independent standard logistic variables (Corollary~\ref{cor:cycle}). In particular
$\sqrt L\,\mathrm{artanh}(u_\infty)\Rightarrow\mathcal N(0,\pi^2/6)$, and $\pi^2/6$ is the expected quadratic variation
of the charge per vertex, on every even cycle and on $\mathbb{Z}$ (Proposition~\ref{prop:line}).
Section~\ref{sec:num} reports simulations: the terminal law on cycles, random limits on $C_4$ and $K_4$ whose dimension
matches the kernel, and the slow decay of the alternating envelope on long cycles that the charges predict for the
line.

\section{The charge martingales}
Let $G=(V,E)$ be a finite graph without loops or isolated vertices, with maximum degree $d$. Each vertex carries an
independent rate-one Poisson clock; when the clock of $v$ rings, one edge $e\ni v$ is chosen with probability $N_e/S_v$,
where $S_v=\sum_{f\ni v}N_f$, and $N_e$ increases by one. All counts start at $1$. Constants $C,C',C'',c$ are finite,
random where stated, and may differ between proofs. Write $X_e(t)=N_e(t)/t$, let $B$ be
the unsigned incidence matrix, $B_{ve}=1$ when $v\in e$, and let $H^{(p)}(n)=\sum_{k=1}^nk^{-p}$, $H=H^{(1)}$.

\begin{theorem}\label{thm:mart}
Let $\sigma\in\mathbb{R}^E$ satisfy $\sum_{e\ni v}\sigma_e=0$ for every $v$. Then
\[ M_\sigma(t)=\sum_{e\in E}\sigma_e\,H(N_e(t)-1) \]
is a martingale with $M_\sigma(0)=0$, jumps bounded by $\max_e|\sigma_e|$, predictable quadratic variation
\[ \langle M_\sigma\rangle_t=\int_0^t\sum_{v\in V}\frac{1}{S_v(s)}\sum_{e\ni v}\frac{\sigma_e^2}{N_e(s)}\,ds , \]
and quadratic variation $[M_\sigma]_t=\sum_e\sigma_e^2H^{(2)}(N_e(t)-1)\le\zeta(2)\sum_e\sigma_e^2$. In particular
$M_\sigma$ is bounded in $L^2$.
\end{theorem}
\begin{proof}
When $N_e$ moves from $n$ to $n+1$, $H(N_e-1)$ increases by $1/n$. A ring at $v$ therefore changes $M_\sigma$ by
$\sigma_e/N_e$ with probability $N_e/S_v$, for each $e\ni v$. The conditional mean of the change is
$S_v^{-1}\sum_{e\ni v}\sigma_e=0$, and its conditional second moment is $S_v^{-1}\sum_{e\ni v}\sigma_e^2/N_e$. On $[0,T]$
the number of rings is Poisson with mean $|V|T$ and every jump is at most $\max_e|\sigma_e|$, so the compensated sums are
square-integrable martingales. The jump at the passage of $N_e$ from $n$ to $n+1$ is $\sigma_e/n$, and each value $n$ is
passed at most once, which gives the formula and the bound for $[M_\sigma]_t$; hence
$\mathbb{E}M_\sigma(t)^2=\mathbb{E}[M_\sigma]_t\le\zeta(2)\sum_e\sigma_e^2$.
\end{proof}

For $\sigma\notin\ker B$ the same computation gives the drift $\sum_v(B\sigma)_v/S_v$, which is not zero in general.
The space $\ker B$ has dimension $|E|-|V|+\beta$, with $\beta$ the number of bipartite components. It is nonzero exactly
when some component contains an even cycle or two odd cycles. Summing the vertex constraints gives $\sum_e\sigma_e=0$
for $\sigma\in\ker B$, and $H(n-1)=\psi(n)+\gamma$, where $\psi$ is the digamma function and $\gamma$ is Euler's constant, so
\begin{equation}\label{eq:link}
\sum_e\sigma_e\log X_e(t)=M_\sigma(t)+\sum_e\sigma_e\bigl(\log N_e(t)-\psi(N_e(t))\bigr),
\end{equation}
with $\log n-\psi(n)=\tfrac1{2n}+O(n^{-2})$. The charges $\sum_e\sigma_e\log X_e$ are therefore martingales up to a term
that vanishes as every count grows. In the mean-field equation $\dot x_e=x_e\sum_{v\in e}S_v(x)^{-1}-x_e$ (time
$\log t$) they are exactly conserved:
$\frac{d}{d\tau}\sum_e\sigma_e\log x_e=\sum_v S_v(x)^{-1}\sum_{e\ni v}\sigma_e-\sum_e\sigma_e=0$.

\begin{lemma}\label{lem:growth}
For an edge $e=(u,v)$ let $a_e(t)=\int_0^t\bigl(S_u(s)^{-1}+S_v(s)^{-1}\bigr)ds$ be its internal time. Then
$D_e(t)=H(N_e(t)-1)-a_e(t)$ converges almost surely and in $L^2$. Almost surely
$\liminf_{t\to\infty}S_v(t)/t\ge1$ for every $v$, and for every $\lambda<2/(d+1)$ and every edge $e$,
$\liminf_{t\to\infty}N_e(t)/t^{\lambda}>0$; in particular every count tends to infinity.
\end{lemma}
\begin{proof}
The count $N_e$ increases at rate $N_e(S_u^{-1}+S_v^{-1})\le2$, and each increase from $n$ to $n+1$ adds $1/n$ to
$H(N_e-1)$, so $D_e$ is a martingale. Its quadratic variation is $H^{(2)}(N_e(t)-1)\le\zeta(2)$, so it is bounded in
$L^2$ and converges almost surely and in $L^2$. Let $F_w(t)$ be the number of rings of $w$ by time $t$. Every ring of
$v$ adds one to $S_v$, so $S_v\ge F_v$ and $\liminf S_v/t\ge1$. A ring of $w$ adds to $S_v$ only if $w=v$ or $w\sim v$,
and then by one, so $S_v(t)\le d+\sum_{w\in\{v\}\cup N(v)}F_w(t)$. Since $F_w(t)/t\to1$ almost surely, for $\varepsilon>0$
there is an almost surely finite $T$ with $S_w(t)\le(d+1)(1+\varepsilon)t$ for all $w$ and $t\ge T$. Then
$a_e(t)\ge\lambda'\log t-C$ for $t\ge T$, with $\lambda'=2/((d+1)(1+\varepsilon))$ and $C$ finite. Since $D_e$
converges, $H(N_e(t)-1)\ge\lambda'\log t-C'$, so $N_e\to\infty$ and, as $H(n-1)-\log n\to\gamma$,
$\log N_e(t)\ge\lambda'\log t-C''$. Taking $\varepsilon$ small gives the claim.
\end{proof}

For $\sigma\in\ker B$ the internal times cancel,
\[ \sum_e\sigma_ea_e(t)=\int_0^t\sum_v\frac{(B\sigma)_v}{S_v(s)}\,ds=0, \]
so that $M_\sigma=\sum_e\sigma_eD_e$.

\begin{theorem}\label{thm:law}
Let $\sigma\in\ker B$. Then $M_\sigma(t)$ converges almost surely and in $L^2$ to a limit $M_{\sigma,\infty}$, and
\[ \mathbb{E}\,e^{i\theta M_{\sigma,\infty}}=\prod_{e\in E}\Gamma(1+i\theta\sigma_e),\qquad\theta\in\mathbb{R}. \]
Equivalently, $M_{\sigma,\infty}$ has the law of $\sum_e\sigma_e\log W_e$, where the $W_e$ are independent standard
exponential variables; since $M_\sigma$ is linear in $\sigma$, the vector $(M_{\sigma,\infty})_{\sigma\in\ker B}$ has the
law of $(\sum_e\sigma_e\log W_e)_{\sigma\in\ker B}$. Moreover $\mathbb{E}M_{\sigma,\infty}^2=
\mathbb{E}\langle M_\sigma\rangle_\infty=\zeta(2)\sum_e\sigma_e^2$ and
$\mathrm{Var}\,\langle M_\sigma\rangle_\infty=\zeta(4)\sum_e\sigma_e^4$.
\end{theorem}
\begin{proof}
Convergence follows from Theorem~\ref{thm:mart}. For real $\theta$ and positive integers $n_e$ put
\[ P_\theta(n)=\prod_e\frac{\Gamma(n_e+i\theta\sigma_e)}{\Gamma(n_e)\,\Gamma(1+i\theta\sigma_e)} . \]
An increase of $n_e$ multiplies $P_\theta$ by $1+i\theta\sigma_e/n_e$, so the generator of the urn applied to $P_\theta$
is $P_\theta(n)\sum_v S_v^{-1}\sum_{e\ni v}n_e\cdot i\theta\sigma_e/n_e=i\theta P_\theta(n)\sum_v(B\sigma)_v/S_v=0$. Since
$|\Gamma(x+iy)|\le\Gamma(x)$ for $x>0$, $|P_\theta|\le\prod_e|\Gamma(1+i\theta\sigma_e)|^{-1}$. A bounded function with
zero generator of a jump process whose total jump rate is $|V|$ is a martingale, so $\mathbb{E}P_\theta(N(t))=1$. As
$n\to\infty$, $\Gamma(n+z)/(\Gamma(n)n^z)\to1$ for fixed complex $z$, and every count tends to infinity
(Lemma~\ref{lem:growth}), so $P_\theta(N(t))\prod_e\Gamma(1+i\theta\sigma_e)-\exp(i\theta\sum_e\sigma_e\log N_e(t))\to0$.
Since $\sum_e\sigma_e=0$, \eqref{eq:link} gives $\sum_e\sigma_e\log N_e(t)-M_\sigma(t)\to0$. By dominated convergence,
$\mathbb{E}e^{i\theta M_{\sigma,\infty}}\prod_e\Gamma(1+i\theta\sigma_e)^{-1}=1$. For a standard exponential $W$,
$\mathbb{E}W^{i\theta}=\Gamma(1+i\theta)$, which gives the second form, and the joint law follows by the Cram\'er--Wold
device.

Since every count tends to infinity, $[M_\sigma]_\infty=\zeta(2)\sum_e\sigma_e^2$, and
$\mathbb{E}M_{\sigma,\infty}^2=\mathbb{E}[M_\sigma]_\infty=\mathbb{E}\langle M_\sigma\rangle_\infty$ by monotone
convergence. The process $R=[M_\sigma]-\langle M_\sigma\rangle$ is a martingale whose jumps are the squared jumps of
$M_\sigma$, and $\langle M_\sigma\rangle$ is continuous, so
$[R]_t=\sum_e\sigma_e^4H^{(4)}(N_e(t)-1)\le\zeta(4)\sum_e\sigma_e^4$. Hence $R$ is bounded in $L^2$, and as
$[M_\sigma]_\infty$ is deterministic,
$\mathrm{Var}\,\langle M_\sigma\rangle_\infty=\mathbb{E}R_\infty^2=\mathbb{E}[R]_\infty=\zeta(4)\sum_e\sigma_e^4$.
\end{proof}

The law in Theorem~\ref{thm:law} is the one obtained by letting each count grow as an independent Yule process run on
its internal time $a_e$: such a process satisfies $\log N_e-a_e\to\log W_e$, and the internal times cancel in every
charge.

\begin{remark}\label{rem:fibre}
Suppose that $X(t)$ converges to a point $x_\infty$ with positive coordinates and deterministic vertex sums
$Bx_\infty=b$. By \eqref{eq:link}, $\sum_e\sigma_e\log x_{\infty,e}=M_{\sigma,\infty}$ for every $\sigma\in\ker B$. These
values determine $x_\infty$: on $P=\{x\in(0,\infty)^E:Bx=b\}$, assumed nonempty, the map $x\mapsto\pi(\log x)$, with $\pi$
the orthogonal projection onto $\ker B$, is a bijection onto $\ker B$. Indeed, for $m\in\ker B$ the strictly convex
function $\sum_ex_e(\log x_e-1)-\langle m,x\rangle$ has a unique minimiser on the compact closure of $P$; it lies in $P$,
because moving from a boundary point towards a point of $P$ decreases the function at infinite rate, and the first-order
condition there is $\pi(\log x)=m$; injectivity follows from $\sum_e(x_e-y_e)(\log x_e-\log y_e)>0$ for distinct
$x,y\in P$. Since $(\ker B)^\perp$ is the range of $B^{\top}$, Theorem~\ref{thm:law} then says that $x_\infty$ has the
law of the unique point of $P$ of the form $x_e=W_e\,e^{y_u+y_v}$, $e=(u,v)$, with vertex potentials $y$: the scaling of
independent exponential edge weights to the vertex sums $b$.
\end{remark}

\section{Convergence on the even cycle}
On the even cycle $C_L$ (edges indexed mod $L$, vertex $v$ between edges $v-1$ and $v$) the kernel is spanned by
$\sigma_e=(-1)^e$, and since $S_v=N_{v-1}+N_v$,
\[ \frac{1}{S_v}\Bigl(\frac1{N_{v-1}}+\frac1{N_v}\Bigr)=\frac{1}{N_{v-1}N_v},\qquad
   \langle M\rangle_t=\int_0^t\sum_{v}\frac{ds}{N_{v-1}(s)N_v(s)} ,\]
where $M=M_\sigma$. Let $L\ge4$ be even. With $K_e$ the compensated count of $e=(u,v)$, and since $S_w(N(t))=tS_w(X(t))$,
\[ dX_e(t)=\frac{F_e(X(t))}{t}\,dt+\frac{dK_e(t)}{t},\qquad F_e(x)=x_e\,\partial_e\mathcal{L}(x), \]
so in the time $\tau=\log t$ the normalised counts follow, up to the martingale term, the field
\begin{equation}\label{eq:ode}
\frac{dx_e}{d\tau}=x_e\,\partial_e\mathcal{L}(x),\qquad
\mathcal{L}(x)=\sum_{v}\log S_v(x)-\sum_e x_e,
\end{equation}
with $\partial_e\mathcal{L}=S_u^{-1}+S_v^{-1}-1$ for $e=(u,v)$. A direct computation \cite{HHKR} gives
\[ \frac{d}{d\tau}\mathcal{L}=\sum_e x_e(\partial_e\mathcal{L})^2\ge0,\qquad \frac{d}{d\tau}\sum_e x_e=L-\sum_e x_e. \]

\begin{lemma}\label{lem:eq}
The equilibria of \eqref{eq:ode} in $\{x\in[0,\infty)^E:\sum_e x_e=L,\ S_v(x)>0\ \forall v\}$ are
\begin{enumerate}
\item the segment $\Sigma=\{x^{(u)}:u\in[-1,1]\}$, $x^{(u)}_e=1+(-1)^eu$, on which $S_v\equiv2$;
\item finitely many points, each with an edge $e$ such that $x_e=0$ and $\partial_e\mathcal{L}(x)\ge 1/(2L)$.
\end{enumerate}
\end{lemma}
\begin{proof}
Let $x$ be an equilibrium. A vertex with both edges at zero has $S_v=0$, which is excluded. If every edge is positive,
$1/S_u+1/S_v=1$ along every edge of the cycle, so $S$ takes two values $s,s'$ alternately; summing,
$\sum_v S_v=2\sum_e x_e=2L$ gives $s+s'=4$, hence $s=s'=2$ and $x\in\Sigma$. Otherwise the positive edges form paths
$w_0w_1\cdots w_k$ separated by single zero edges, and along such a path $S(w_i)$ alternates between $s$ and
$s'=s/(s-1)$. Since the edges outside the path at $w_0$ and $w_k$ are zero, $x_1=S(w_0)=s$, and solving along the path
gives $x_{2i}=id$, $x_{2i+1}=s-id$ with $d=s'-s$. If $k$ is odd the end condition $x_k=S(w_k)=s'$ forces $d=0$, hence
$k=1$ and $x_1=2$. If $k$ is even it forces $d=2s/k$, and $1/s+1/s'=1$ gives $s=(2k+2)/(k+2)$; all $x_j$ are then
positive. So every block is either a single edge of weight $2$ with $S=2$ at both ends, or an even block of length $k$
with $S=(2k+2)/(k+2)$ at both ends, and the configuration is determined by the block lengths. If all blocks have length
one, $x$ alternates $2,0$ and lies in $\Sigma$ ($u=\pm1$). Otherwise some even block of length $k\le L-1$ meets a zero edge
$e$ whose other end has $S\in\{2,(2k'+2)/(k'+2)\}$, so
$\partial_e\mathcal{L}\ge \frac{k+2}{2k+2}+\frac12-1=\frac{1}{2k+2}\ge\frac1{2L}$, the smallest value occurring when
the other end has $S=2$.
\end{proof}

\begin{theorem}\label{thm:conv}
On $C_L$, $L$ even, almost surely $X(t)\to x^{(u_\infty)}$ with $|u_\infty|<1$ and
\[ L\,\mathrm{artanh}(u_\infty)=M_\infty:=\lim_{t\to\infty}\sum_e(-1)^eH(N_e(t)-1). \]
\end{theorem}
\begin{proof}
By Lemma~\ref{lem:growth} and since $\sum_eN_e(t)-L$ is a Poisson process of rate $L$, $X(t)$ eventually stays in the
compact set $K=\{x\in[0,\infty)^E:\sum x_e\le 2L,\ S_v(x)\ge1/2\ \forall v\}$, which is forward invariant for
\eqref{eq:ode} and on which the field is Lipschitz. The total jump rate is $L$, so by Doob's inequality the martingale
term over a window satisfies
$\mathbb{E}\sup_{h\le T}|\int_{e^\tau}^{e^{\tau+h}}t^{-1}dK_e(t)|^2\le4\int_{e^\tau}^{e^{\tau+T}}Lt^{-2}\,dt\le4Le^{-\tau}$,
and the Borel--Cantelli lemma along integer $\tau$ shows that it tends to zero almost surely, uniformly over windows of
length $T$. The jumps of $X$ are at most $1/t$. By Gronwall's lemma the linear interpolation of $\tau\mapsto X(e^\tau)$ is
then a precompact asymptotic pseudotrajectory of the flow of \eqref{eq:ode} on $K$ in the sense of
\cite[Section~3]{Ben99}, and its limit set is internally chain transitive \cite[Theorem~5.7]{Ben99}. The function
$-\mathcal{L}$ is strictly decreasing along every nonstationary orbit, and by Lemma~\ref{lem:eq} it takes finitely many
values on the equilibria ($\mathcal{L}\equiv L\log2-L$ on $\Sigma$). By \cite[Proposition~6.4]{Ben99} the limit set
consists of equilibria, and being connected it lies in $\Sigma$ or is a single point of type (2).

Type (2): let $x^*$ be such a point, with $x^*_e=0$ and $\partial_e\mathcal{L}(x^*)\ge1/(2L)$, and choose a
neighbourhood $U$ of $x^*$ in $K$ on which $x_e\le1$ and $\partial_e\mathcal{L}\ge1/(4L)$. On the event
$A_m=\{X(t)\in U\text{ for all }t\ge m\}$, $m\in\mathbb{N}$, the internal time of $e$ satisfies
$a_e(t)-a_e(m)=\int_m^t(1+\partial_e\mathcal{L}(X(s)))\,s^{-1}ds\ge(1+\delta)\log(t/m)$ with $\delta=1/(4L)$. By
Lemma~\ref{lem:growth}, $\log N_e(t)-a_e(t)$ converges, so $X_e(t)\ge ct^{\delta}\to\infty$ on $A_m$, which contradicts
$x_e\le1$ on $U$. Hence $\mathbb{P}(A_m)=0$ for every $m$, and $X$ does not converge to $x^*$. As there are finitely many
such points, the limit set lies in $\Sigma$.

By \eqref{eq:link} and Theorem~\ref{thm:mart}, $\Phi(X(t)):=\sum_e(-1)^e\log X_e(t)\to M_\infty$, finite almost surely.
Along a sequence of times with $X\to x^{(1)}$ every even coordinate tends to $2$ and every odd coordinate to $0$, so
$\Phi\to+\infty$; similarly $\Phi\to-\infty$ near $x^{(-1)}$. So neither endpoint is a limit point, and the limit set is a
compact subset of $\{x^{(u)}:|u|<1\}$, where $\Phi$ is continuous and
$\Phi(x^{(u)})=\tfrac L2\log\frac{1+u}{1-u}=L\,\mathrm{artanh}(u)$ is strictly increasing in $u$. Every limit point
$x^{(u)}$ therefore has $L\,\mathrm{artanh}(u)=M_\infty$, and the limit set is the single point $x^{(u_\infty)}$.
\end{proof}

\begin{corollary}\label{cor:cycle}
On $C_L$, $L\ge4$ even, $u_\infty=\tanh(M_\infty/L)$, where $M_\infty$ has the law of $\Lambda_1+\dots+\Lambda_{L/2}$
with independent standard logistic variables $\Lambda_j$ (density $e^{-s}/(1+e^{-s})^2$), that is,
\[ \mathbb{E}\,e^{i\theta M_\infty}=\Bigl(\frac{\pi\theta}{\sinh\pi\theta}\Bigr)^{L/2}. \]
In particular $\mathbb{E}M_\infty^2=L\pi^2/6$, $\mathbb{E}\bigl(\langle M\rangle_\infty/L-\pi^2/6\bigr)^2=\pi^4/(90L)$, and as
$L\to\infty$ through even values, both $\sqrt L\,\mathrm{artanh}(u_\infty)$ and $\sqrt L\,u_\infty$ converge in law to
$\mathcal N(0,\pi^2/6)$.
\end{corollary}
\begin{proof}
By Theorem~\ref{thm:law} with $\sigma_e=(-1)^e$, $M_\infty$ has the law of $\sum_{j}(\log W_{2j}-\log W_{2j+1})$, and
$\mathbb{E}e^{i\theta(\log W-\log W')}=\Gamma(1+i\theta)\Gamma(1-i\theta)=\pi\theta/\sinh(\pi\theta)$ for independent
standard exponential $W,W'$, the characteristic function of the standard logistic law. Theorem~\ref{thm:conv} gives
$u_\infty=\tanh(M_\infty/L)$, and Theorem~\ref{thm:law} gives the two moments. The logistic law has variance $\pi^2/3$,
so the central limit theorem gives $M_\infty/\sqrt L\Rightarrow\mathcal N(0,\pi^2/6)$, and
$\sqrt L\,u_\infty-M_\infty/\sqrt L\to0$ in probability since $\tanh(x)/x\to1$ as $x\to0$ and $M_\infty/L\to0$.
\end{proof}

On $C_4$, with $p=(1+u_\infty)/2$, the density of $M_\infty$ is $f(s)=(s\coth(s/2)-2)/(4\sinh^2(s/2))$, the density of
$p$ is $2f(2\log\frac{p}{1-p})/(p(1-p))$, which behaves like $4(1-p)\log\frac1{1-p}$ as $p\to1$, and
$\mathrm{Var}\,p=\pi^2/32-1/4=0.05843\dots$. The limit law on $C_4$ is therefore not a beta law.

\section{The infinite line}\label{sec:line}
On $\mathbb{Z}$, with vertex $v$ between edges $v-1$ and $v$, the process is built from the same rate-one clocks and
choices. A vertex $w$ that does not ring during $[0,T]$ separates the line on that interval: the edges to the left of $w$
change only at rings to the left of $w$, which depend only on those edges, and similarly on the right. Each vertex is
silent on $[0,T]$ with probability $e^{-T}$, independently, so almost surely the line splits into finite blocks, each of
which evolves on $[0,T]$ as a finite urn. This defines the process for all times, and its law is invariant under
translations.

\begin{proposition}\label{prop:line}
For the critical WARM on $\mathbb{Z}$ with unit initial counts every count tends to infinity almost surely, and
\[ c:=\mathbb{E}\int_0^\infty\frac{dt}{N_{-1}(t)N_0(t)}=\zeta(2)=\frac{\pi^2}{6}. \]
The same expectation, with $N_{v-1}N_v$ in place of $N_{-1}N_0$, equals $\pi^2/6$ for every vertex $v$ of every even cycle
$C_L$, $L\ge4$.
\end{proposition}
\begin{proof}
The proof of Lemma~\ref{lem:growth} uses only the clocks of finitely many vertices near the edge, so it applies on
$\mathbb{Z}$ with $d=2$ and every count tends to infinity. The count $N_0$ increases at rate $N_0(S_0^{-1}+S_1^{-1})\le2$,
and each increase from $n$ to $n+1$ adds $n^{-2}$ to $A(t)=H^{(2)}(N_0(t)-1)$, so
$\mathbb{E}A(T)=\int_0^T\mathbb{E}\bigl[(N_0S_0)^{-1}+(N_0S_1)^{-1}\bigr]dt$. By translation invariance
$\mathbb{E}[(N_0S_1)^{-1}]=\mathbb{E}[(N_{-1}S_0)^{-1}]$, and $(N_0^{-1}+N_{-1}^{-1})/S_0=1/(N_{-1}N_0)$, so
$\mathbb{E}\int_0^T\frac{dt}{N_{-1}N_0}=\mathbb{E}H^{(2)}(N_0(T)-1)$, which increases to $\zeta(2)$ by bounded
convergence. On $C_L$, $\mathbb{E}\langle M\rangle_\infty=L\zeta(2)$ by Theorem~\ref{thm:law}, and the $L$ vertex terms
have the same expectation by rotation invariance.
\end{proof}

On $\mathbb{Z}$ the kernel is replaced by a local conservation law. With $Y_e=(-1)^eH(N_e-1)$ and $J_v=(-1)^v/S_v$,
the computation in Theorem~\ref{thm:mart} gives
\[ dY_e=-(J_{e+1}-J_e)\,dt+dM_e , \]
with $M_e$ a martingale. Linearising \eqref{eq:ode} at $x\equiv1$ gives the decay rate $(1+\cos k)/2$ for the mode
$e^{ike}$, which vanishes at $k=\pi$. Writing $x_e=1+(-1)^eu(e)+w(e)$ with $u$ slowly varying, the field $w$ relaxes to
$(u^2)''/8$ and, formally, $u$ obeys
\[ \partial_\tau u=\tfrac14(1-u^2)\,\partial_e^2u ,\qquad\text{equivalently}\qquad \partial_\tau v=\tfrac14\,\partial_e^2\tanh v,\quad v=\mathrm{artanh}\,u, \]
which conserves $\int v$, the continuum form of the charge. Started from the noise of the early stage, this heat flow
in $\tau=\log t$ predicts that the local alternating amplitude decays like $(\log t)^{-1/4}$.

\section{Numerical checks}\label{sec:num}
All simulations run the process exactly (rate-one vertex clocks, unit initial counts).

\emph{Terminal law on cycles.} The program \texttt{code/exact\_law\_check.py} runs independent copies of the jump chain
of the urn on $C_L$ and compares the charge after $K$ rings with the law of Corollary~\ref{cor:cycle}. With $K=4000$ and
$20000$ copies on $C_4$, the variance of the charge is $6.63$ (exact $2\pi^2/3=6.58$, standard error about $0.07$), its
excess kurtosis $0.59$ (exact $12/(5L)=0.6$), the Kolmogorov--Smirnov distance to the exact law $0.005$ (the $5\%$ level
is about $0.010$), and $\mathrm{Var}\,p=0.0587$ for $p=(1+\tanh(M/4))/2$ (exact $0.0584$). With $K=10^5$ and $4000$
copies on $C_{50}$ the variance is $81.9$ (exact $82.2$, standard error $1.8$) and the Kolmogorov--Smirnov distance
$0.008$ (the $5\%$ level is about $0.022$).

\emph{Even cycles in continuous time.} The program \texttt{code/warm2.c} runs $R$ independent copies to time $T$ (each
run stops at the first ring after $T$) and records $M(T)$, the ring-sampled bracket and $\mathrm{artanh}(u)$. The
ring-sampled bracket is the sum over rings of $1/(N_{v-1}N_v)$ for the ringing vertex $v$ before the ring; its mean is
$\mathbb{E}\langle M\rangle_T$. The value $u$ is read from $\sum_e(-1)^e\log X_e(T)=L\,\mathrm{artanh}(u)$.
\begin{center}
\begin{tabular}{rrrrrrr}
\toprule
$L$ & $T$ & $R$ & mean $M$ (s.e.) & $\mathrm{Var}\,M$ & mean bracket & $L\,\mathbb{E}\,\mathrm{artanh}(u)^2$\\
\midrule
100 & $10^3$ & 400 & $0.30$ $(0.64)$ & 161.6 & 164.3 & 1.61\\
400 & $10^3$ & 200 & $-2.86$ $(1.76)$ & 620 & 657 & 1.56\\
100 & $10^4$ & 200 & -- & 154.9 & 164.7 & --\\
\bottomrule
\end{tabular}
\end{center}
The terminal values are $\mathbb{E}M_\infty^2=\mathbb{E}\langle M\rangle_\infty=L\pi^2/6$, that is $164.5$ and $658.0$,
and $L\,\mathbb{E}\,\mathrm{artanh}(u_\infty)^2=\pi^2/6=1.645$ for every $L$. The bracket per vertex, $1.64$ at both sizes,
agrees with $\pi^2/6$; the sample variances have standard errors of about $10\%$ of their size at these $R$. The difference
$\sum_e(-1)^e\log X_e(T)-M(T)$ in \eqref{eq:link} has root mean square $0.007$ at $T=10^4$, $L=100$. At $T=10^3$ the counts
on $C_{100}$ are not yet a pure alternation, which lowers the last column: the slow modes of \eqref{eq:ode} near
wavenumber $\pi$ decay like $t^{-\pi^2/L^2}$.

\emph{Small graphs.} On $K_4$, where $\dim\ker B=2$, the standard deviation of each $X_e$ over $2000$ runs stays near
$0.35$ between $T=10^3$ and $T=10^5$, while $\max_v|S_v/t-2|$ falls from $0.25$ to $0.08$: the limit is a random point
of a two-dimensional set of equilibria, as Theorem~\ref{thm:law} and Remark~\ref{rem:fibre} describe. On the triangle
($\ker B=0$) the spread decreases, more slowly than $t^{-1/4}$, towards $(1,1,1)$.

\emph{The line.} On a cycle of length $1001$ run to $T=10^5$ the standard deviation of $X_e$ falls from $0.57$ at
$t=12$ to $0.39$ at $t=8.8\times10^4$, and $(\log t)^{1/4}$ times it stays between $0.67$ and $0.72$, rising slowly (a
fitted exponent of $-0.18$ in $\log t$ over this range, which a shift of $\tau$ by about $3$ or the saturation of
$\tanh$ would account for); the correlation of neighbouring edges tends to $-0.92$. These runs are consistent with
Kleptsyn's conjecture, with convergence at the logarithmic rate the charges predict.

\section*{Acknowledgements and funding}
The research presented in this paper was supported by the European Research Council (ERC) under the European Union's
Horizon Europe research and innovation programme (grant agreement No.~101041711), by the Simons Foundation as part of
the Collaboration on the Mathematical and Scientific Foundations of Deep Learning, by Heights Labs, by Convex Nexus
Capital, by the Israel Science Foundation (grant number 2258/19), and by the Israel Science Foundation (ISF Grant
4101/25). Computational resources were provided in part by Google Cloud through its education credits programme (credit
530550103).

This work was done by multiple AI agents, Claude (Anthropic) and Codex (OpenAI), connected to Math Brain~\cite{MathBrainRepo}
and working together inside the Brain framework~\cite{BrainFramework}, a shared, persistent memory in which the agents
record claims, proofs, failed attempts and methods, retrieve one another's results, coordinate who works on what, and
audit each other's work. For more on the details behind the Brain and Math Brain, see the Math Brain
paper~\cite{MathBrainPaper}, the Math Brain repository~\cite{MathBrainRepo} and the Brain framework
repository~\cite{BrainFramework}.

\section*{Declaration of competing interest}
The author holds equity in Heights Labs and in Convex Nexus Capital and has a management role at Convex Nexus Capital.
These entities are included in the funding declaration above.

\section*{Declaration of generative AI and AI-assisted technologies}
The paper comes out of Math Brain, an ongoing project on the use of AI systems for research mathematics. In November
2026 Claude Opus~5.5 (Anthropic), working in Claude Code, found the charge martingales of Theorem~\ref{thm:mart}, proved
the convergence on the even cycle, wrote and ran the simulations of Section~\ref{sec:num}, and wrote the paper. A Codex
agent (OpenAI), refereeing the paper, found the terminal law of Theorem~\ref{thm:law} and its proof, the identities for
$[M_\sigma]$ and $\mathrm{Var}\,\langle M_\sigma\rangle_\infty$, the identification $c=\zeta(2)$ in
Proposition~\ref{prop:line} and the convex proof in Remark~\ref{rem:fibre}. The author directed the work and is
responsible for the content of the paper and for decisions concerning its dissemination.

\section*{Data availability}
No data were used. The code is available at\\ \url{https://github.com/yspennstate/graph-urn-conserved-charges}.

The paper is part of the Math Brain project, whose public repository is
\url{https://github.com/yspennstate/math-brain-public}~\cite{MathBrainRepo}.

\appendix
\section{Notation}
\begin{center}\footnotesize
\begin{tabular}{l p{6.3cm} l}
\toprule
Symbol & Meaning & First defined\\
\midrule
$G=(V,E)$, $d$ & finite graph, its maximum degree & Section 2\\
$N_e$, $X_e$ & count of edge $e$, normalised count $N_e(t)/t$ & Section 2\\
$S_v$ & vertex sum of counts $\sum_{e\ni v}N_e$ & Section 2\\
$B$ & unsigned vertex--edge incidence matrix & Section 2\\
$H^{(p)}$, $H$ & $H^{(p)}(n)=\sum_{k=1}^nk^{-p}$, harmonic number $H=H^{(1)}$ & Section 2\\
$\sigma\in\ker B$ & charge vector & Thm.~\ref{thm:mart}\\
$M_\sigma$, $M_{\sigma,\infty}$ & charge martingale, its terminal value & Thms.~\ref{thm:mart}, \ref{thm:law}\\
$\langle M_\sigma\rangle$, $[M_\sigma]$ & predictable and optional quadratic variation & Thm.~\ref{thm:mart}\\
$\psi$, $\gamma$ & digamma function, Euler's constant & \eqref{eq:link}\\
$F_w(t)$ & number of rings of the clock of $w$ by time $t$ & Lem.~\ref{lem:growth}\\
$a_e(t)$ & internal time of edge $e$ & Lem.~\ref{lem:growth}\\
$D_e$ & edge martingale $H(N_e-1)-a_e$ & Lem.~\ref{lem:growth}\\
$W_e$ & independent standard exponential variables & Thm.~\ref{thm:law}\\
$P_\theta(n)$ & Gamma-ratio martingale & Thm.~\ref{thm:law}\\
$R$ & compensated quadratic variation $[M_\sigma]-\langle M_\sigma\rangle$ & Thm.~\ref{thm:law}\\
$P$, $b$, $\pi$ & fibre $\{x>0:Bx=b\}$, its vertex sums, orthogonal projection onto $\ker B$ & Rem.~\ref{rem:fibre}\\
$C_L$ & cycle with $L$ edges, $L$ even & Section 3\\
$\tau$ & time $\log t$ of the mean-field equation & \eqref{eq:ode}\\
$\mathcal{L}$ & Lyapunov function of \eqref{eq:ode} & \eqref{eq:ode}\\
$K_e$ & compensated count of edge $e$ & Section 3\\
$x^{(u)}$, $\Sigma$ & equilibrium $x^{(u)}_e=1+(-1)^eu$, segment of equilibria & Lem.~\ref{lem:eq}\\
$u_\infty$, $M_\infty$ & limit point on $\Sigma$, terminal charge on $C_L$ & Thm.~\ref{thm:conv}\\
$\Phi$ & $\sum_e(-1)^e\log x_e$, the charge as a function of the normalised counts & Thm.~\ref{thm:conv}\\
$\Lambda_j$ & independent standard logistic variables & Cor.~\ref{cor:cycle}\\
$p$ & limit share $(1+u_\infty)/2$ on $C_4$ & Section 3\\
$c$ & expected bracket per vertex on $\mathbb{Z}$ & Prop.~\ref{prop:line}\\
\bottomrule
\end{tabular}
\end{center}

\bibliographystyle{amsplain}
\begin{thebibliography}{99}
\bibitem{AletOP} G.~Aletti, A.~Collevecchio, I.~Crimaldi, M.~Deijfen, S.~Dereich, A.~Ghiglietti, R.~van der Hofstad,
M.~Holmes, D.~Kious, V.~Kleptsyn, B.~Lodewijks, P.~M\"orters and B.~Schapira, Open problems in stochastic
reinforcement models, in 2023 MATRIX Annals, MATRIX Book Series, Springer, 2025, 477--498.
\bibitem{BBCL} M. Bena\"im, I. Benjamini, J. Chen, Y. Lima, A generalized P\'olya's urn with graph based interactions, Random Structures Algorithms 46 (2015), 614--634, doi:10.1002/rsa.20523.
\bibitem{CL} J. Chen, C. Lucas, A generalized P\'olya's urn with graph based interactions: convergence at linearity, Electron. Commun. Probab. 19 (2014), 1--13, doi:10.1214/ECP.v19-3094.
\bibitem{HHKR} R. van der Hofstad, M. Holmes, A. Kuznetsov, W. Ruszel, Strongly reinforced P\'olya urns with graph-based competition, Ann. Appl. Probab. 26 (2016), 2494--2539, doi:10.1214/16-AAP1153.
\bibitem{Lima} Y. Lima, Graph-based P\'olya's urn: completion of the linear case, Stoch. Dyn. 16 (2016), no.~2, 1660007, doi:10.1142/S0219493716600078.
\bibitem{Ben99} M. Bena\"im, Dynamics of stochastic approximation algorithms, S\'eminaire de Probabilit\'es XXXIII, Lecture Notes in Math. 1709, Springer (1999), 1--68.
\bibitem{MathBrainRepo} Y.~Shmalo, \emph{Math Brain: public snapshot, code and tests}, 2026, \url{https://github.com/yspennstate/math-brain-public}.
\bibitem{MathBrainPaper} Y.~Shmalo, \emph{Math Brain: shared, auditable memory for multi-agent mathematical research}, October 2026, doi:10.5281/zenodo.24150837, \url{https://github.com/yspennstate/math-brain-public/tree/main/papers/math-brain-2026-10-24}.
\bibitem{BrainFramework} Y.~Shmalo, \emph{Brain Framework}, public software repository, 2026, \url{https://github.com/yspennstate/brain-framework-public}.
\end{thebibliography}
\end{document}
